% META: {"id": "1NSI-ARCHOS-ME-004", "chapitre": "1NSI-ARCHITECTURE-OS", "type_objet": "methode", "capacites": ["P-ARCH-03C"], "capacites_codes": ["C5"], "methodes": ["M4"], "mode_creation": "ex_nihilo", "sources_inspiration": [], "status": "needs_review"}
% BEGIN-VERIFY
% def triplets(permissions):
%     return [permissions[0:3], permissions[3:6], permissions[6:9]]
% p = "rw-r-xr--"
% assert triplets(p) == ["rw-", "r-x", "r--"]
% def autorise(permissions, categorie, action):
%     indices = {"proprietaire": 0, "groupe": 1, "autres": 2}
%     positions = {"lire": 0, "ecrire": 1, "executer": 2}
%     bloc = triplets(permissions)[indices[categorie]]
%     return bloc[positions[action]] != "-"
% assert autorise(p, "proprietaire", "ecrire") is True
% assert autorise(p, "groupe", "ecrire") is False
% assert autorise(p, "groupe", "executer") is True
% assert autorise(p, "autres", "lire") is True
% assert autorise(p, "autres", "executer") is False
% END-VERIFY
\begin{fichemethode}{M4}{Lire des permissions de fichier sans se tromper de triplet}

\textbf{Quand l'utiliser ?} Devant une écriture comme \lstinline{rw-r-xr--}, pour dire qui a le
droit de faire quoi.

\textbf{Pas à pas :}
\begin{enumerate}
  \item \textbf{Découpe en trois blocs de trois}, toujours dans le même ordre :
        caractères $1$ à $3$ le \textbf{propriétaire}, $4$ à $6$ le \textbf{groupe},
        $7$ à $9$ les \textbf{autres}. Écris ce découpage avant de répondre.
  \item \textbf{Dans chaque bloc, l'ordre est fixe} : \lstinline{r} lire, \lstinline{w} écrire,
        \lstinline{x} exécuter. Un tiret signifie que le droit est refusé, à cette position-là
        et à aucune autre.
  \item \textbf{Repère d'abord la catégorie demandée}, puis seulement la position dans son
        bloc. C'est l'inversion de ces deux étapes qui produit la plupart des erreurs.
  \item \textbf{Relis le tiret} : un tiret en deuxième position d'un bloc refuse l'écriture,
        pas la lecture.
\end{enumerate}

\textbf{Exemple rédigé.} Pour \lstinline{rw-r-xr--} : le propriétaire a \lstinline{rw-}, donc
il lit et écrit mais n'exécute pas. Le groupe a \lstinline{r-x} : il lit et exécute, mais
\textbf{n'écrit pas} — le tiret est en deuxième position. Les autres ont \lstinline{r--} :
lecture seule.

\textbf{Erreur fréquente.} Lire \lstinline{rw} là où figure \lstinline{r-x}, ou répondre en
lisant le premier triplet — celui du propriétaire — quand la question porte sur le groupe, qui
est le deuxième.

\end{fichemethode}
